\section{Cinemática directa e inversa}
La \textbf{cinemática} de un robot estudia el movimiento de sus partes sin considerar las fuerzas que lo producen. 

Para una \textbf{cadena cinemática abierta simple}, la \textbf{cinemática directa} calcula la posición y orientación del efector final respecto del \textbf{marco de referencia del mundo} a partir de los ángulos o movimientos de sus articulaciones. La \textbf{cinemática inversa} hace lo contrario, calcula qué valores deben tener los ángulos o movimientos de las articulaciones para que el efector final llegue a un punto deseado respecto del \textbf{marco de referencia del mundo}. \todo[inline]{basado de una visión general de IA de google con el prompt -cinemática directa e inversa de un robot-}

El método D-H es necesario para, de entre otras cosas, resolver el problema de la cinemática directa. El flujo para resolver el problema de la cinemática directa es este:
\begin{enumerate}
    \item Tomar los ángulos o desplazamientos de las articulaciones
    \item Construir todas las matrices de transformación
    \item Ubicar el efector final 
\end{enumerate}